\begin{circuitikz}[
    american,
    voltage dir=RP,
    >=latex,
    %>={Latex[length=5mm, width=2mm]},
    alt={Ladder logic diagram showing a single rung connecting two vertical power rails. On the left side of the rung, a normally open contact labeled StartPB is in parallel with a normally open contact labeled AutoMode. These parallel inputs reconnect before continuing along the rung to drive an output coil labeled GreenLT on the far right.}
    ]

    % Define the number of rungs in the diagram
    \def\numRungs{1}

    % Calculate the length of the vertical power rails dynamically based on number of rungs
    \pgfmathsetmacro{\railLength}{-2 * \numRungs}

    %======================================================================================================
    % Component Grid
    % Spacing = 3x, 2y
    %======================================================================================================
    \foreach \i in {0,...,5} {
        \foreach \j in {0,...,20} { % Change to 0,...,19 if you need 20 individual Y points down to -40
        
        % Calculate spatial positions
        \pgfmathsetmacro{\px}{\i * 3}
        \pgfmathsetmacro{\py}{-1 - (\j * 2)}
        
        % Name coordinate using loop indices: (C-column-row) e.g., (C-0-0), (C-5-10)
        \coordinate (C-\i-\j) at (\px, \py);
        
        % OPTIONAL: Uncomment the 2 lines below to display grid labels while building
        % \fill[red] (C-\i-\j) circle (1.5pt);
        %\node[anchor=south west, scale=0.5, gray] at (C-\i-\j) {(\i,\j)};
        }
    }

    %======================================================================================================
    % Foreground Layer
    % Only contains the vertical power rails
    % Length of the rails is automatic, calculated based on the number of rungs defined above
    %======================================================================================================
    \begin{pgfonlayer}{ft}

        %======================================================================================================
        % Scope sets options for both lines - BakerBlack, very thick
        %======================================================================================================
        \begin{scope}[
            draw=BakerBlack,
            text=BakerBlack,
            color=BakerBlack,
            font=\small,
            ultra thick
        ]
        
            \draw (0,0) -- (0,\railLength);

            \draw (15,0) -- (15,\railLength);

            

        \end{scope}
    \end{pgfonlayer}

    %======================================================================================================
    % Main Layer
    % All actual logic goes here
    %======================================================================================================
    \begin{pgfonlayer}{main}
        
        %======================================================================================================
        % Scope sets options for both lines - BakerBlack for draw and text
        %======================================================================================================
        \begin{scope}[
            draw=BakerBlack,
            text=BakerBlack,
            color=BakerBlack,
            font=\small
        ]

            %=================================================================
            % Rung 1
            %=================================================================
            \draw (C-0-0) to[C,l={StartPB}] (C-1-0) -- (C-4-0) to[esource,l={GreenLT}] (C-5-0);

            \draw ($(C-0-0)+(0.25,0)$) -- ++(0,-2) to[C,l={AutoMode}] ++(2.75,0) -- (C-1-0);

        \end{scope}
    \end{pgfonlayer}

    %=================================================================
    % Background layer
    %=================================================================
    \begin{pgfonlayer}{bg}
        
        %==================================================================
        % Highlight Scope
        % This layer exists for highlighting true logic, the scope sets the
        % highlight options without having to repeat them every \draw
        % draw=UpNorthGreen, line width=10pt, opacity=0.4
        %==================================================================
        \begin{scope}[
            draw=UpNorthGreen,
            line width=10pt,
            opacity=0.4
        ]

            

        \end{scope}
    \end{pgfonlayer}
    
\end{circuitikz}